\section{Implementation details of NOVA and the injector}
\label{app:implementation}

% Question: Where are the settings, tolerances and checks of the current
% configuration recorded?
% Answer: Here, so that the Methods can explain the ideas while every number
% remains available. Each subsection mirrors a subsection of the Methods.
% Evidence: As listed in the comment blocks of sections/02_methods.tex; the
% content-change table proposals/CONTENT_CHANGES_20260924.md.

This appendix collects the settings, tolerances and checks of the current
configuration which are referred to, but not listed, in
Section~\ref{sec:methods}. The subsections follow the order of that section.

\subsection{NOVA}

\paragraph{Support and wavelength grid.}
The support of the model contains 97{,}557 pixel positions per integration,
in 2{,}720 groups, after six groups at the edge of the detector without a
reliable wavelength mapping have been excluded. Individual samples are
removed if their value is not finite, their uncertainty is not positive, or
they carry the \texttt{DO\_NOT\_USE} flag of the JWST data-quality array. The
supports of the two orders share no pixel. The transit is evaluated on a
fine grid of 5{,}454 wavelength nodes.

\paragraph{Spectral basis.}
The 160 top-hat basis functions span 0.62 to 2.84~$\mu\mathrm{m}$; they are
about 0.01~$\mu\mathrm{m}$ wide below 2~$\mu\mathrm{m}$ and about
0.04~$\mu\mathrm{m}$ wide above it. The two coefficients above
2.76~$\mu\mathrm{m}$ have no response support inside the validated
wavelength calibration and are held constant, which leaves 158 free depth
coordinates. The earlier release had an order-difference term, an offset and
a slope in log-wavelength added to order 1 and subtracted from order 2,
which was disabled by bounds of $\pm10^{-12}$. It is now fixed at exactly
zero, because at those small but nonzero bounds the derivatives of the
penalty terms were not exactly those of a model without it.

\paragraph{Orbit and transit model.}
The period, $P=3.73548546$~d, is the value of the public Ahsoka configuration
of \citet{LouieEtAl2025}. The argument of periastron is set to $90^\circ$,
which has no effect for a circular orbit, so the value of zero listed for
the earlier release in Appendix~\ref{app:nova} is equivalent. jaxoplanet
evaluates the occulted flux with a quadrature of order 10, and the average
over each exposure uses 21 Gauss--Legendre samples per integration.

\paragraph{Limb-darkening reference.}
The STAGGER reference uses the 6500~K plane of the grid, the closest complete
plane to the stellar effective temperature of 6550~K, interpolated
bilinearly in $\log g$ and $[\mathrm{Fe}/\mathrm{H}]$ to $\log g=4.149$ and
$[\mathrm{Fe}/\mathrm{H}]=-0.25$. The reference file is fixed and identified
by its content hash. The deviation from theory is penalised with a scale of
0.5 in logit space, and its curvature in wavelength with a scale of
$2.0\,\mu\mathrm{m}^{-2}$. Fine wavelength nodes outside the range of the
theory grid are excluded, and the theory is never extrapolated.

\paragraph{Spatial response.}
The template $\phi_o$ is built from the bright, stable pixels of each order:
those with at least 95\% of their out-of-transit samples valid, a positive
baseline, and an out-of-transit level above the 35th percentile of the
stable pixels of that order. For each, the fractional departure from its own
linear out-of-transit baseline is computed; the template is the median of
these departures at each integration, with any linear trend removed and
scaled so that its deepest point is exactly $-1$. Every pixel of the support
is then fitted with Equation~\ref{eq:qstar-fit}. The time coordinate
$\tau_t$ is measured from the mean out-of-transit time, in units of the
out-of-transit time span. The normalised amplitudes are smoothed across
neighbouring columns with a robust second-difference penalty, at a fixed
rank of the pixel when the pixels of each group are numbered by detector row.
The smoothing strength is the one for which the response smoothed on one of
two alternating folds, stratified across the transit phases, best predicts
the unsmoothed response of the other. The response used in the fit is built
from all integrations with this strength, and it is only accepted if the two
fold estimates agree in their flux distribution and centroid
(Appendix~\ref{app:nova}). Negative values are then set to zero, the
response of each column is combined with the positive transit-correlated
amplitudes measured in the wings of the same column, in proportion to their
amplitude, and the result is renormalised within each group; groups without
a smoothed response use the amplitudes alone.

\paragraph{Background.}
The eight maps are fitted by weighted least squares to the 244{,}657
off-trace background pixels of each integration separately, which gives a
time series of eight coefficients. The seven polynomial maps are
orthonormalised on the off-trace pixels. $G_1$ is the first of the principal
components of this series whose absolute correlation with the mean of
$\phi_o$ over the two orders is at most 0.25. The anchor
$\bm{\alpha}_{\mathrm{off}}$ and its precision $\bm{N}_{\mathrm{off}}$ come
from one joint weighted fit of all sixteen coefficients to every off-trace
pixel of every integration.

\paragraph{Source-curvature factor.}
The factor is built from the background-subtracted flux of the stable trace
pixels of each order, i.e.\ those which are valid in every out-of-transit
integration, and NOVA uses the geometric mean of the factors of the two
orders. Its common form was chosen on the earlier release by comparing its
held-out $\chi^2$ on two folds with that of no correction and of a separate
factor for each order. On each input, two further factors are computed from
two alternating folds of the out-of-transit integrations (175 each for the
real visit). The factor is accepted only if the means of the two fold
factors over the in-transit integrations differ by at most 250~ppm and if
its largest deviation from unity is at most 1000~ppm. For the real visit and
the injected inputs, the two fold factors must also be correlated with a
coefficient of at least 0.95; for the ensemble realisations this correlation
is only reported. That rule was changed after the correlation was found to
reject realisations whose corrections were only 1--15~ppm, since two
negligible corrections of opposite sign are strongly anticorrelated even
when both amplitude limits are met.

\paragraph{Uncertainties.}
The per-order scales are $N_1=3.07$ and $N_2=4.19$. The per-pixel factor
$s_p$ is estimated by two-way cross-validation on alternating folds of the
out-of-transit integrations: a level and a slope in time are fitted to each
pixel on one fold and used to predict the other, and vice versa, and $s_p$
is set by the scatter of the standardised prediction residuals, with a
minimum of one.

\paragraph{Convergence and starts.}
Between successive outer iterations, the Huber weights must change by less
than $10^{-3}$ and the objective by less than $10^{-8}$ in relative terms,
and the spectrum, mapped onto two fixed summary grids inherited from the
earlier release (40 and 96 bins), must change by less than 0.01 and
0.05~ppm root-mean-square. First-order optimality must also be satisfied,
the rank of the linear block must be stable, the solution must be physically
feasible, and a confirmation cycle with the recomputed final weights must
reproduce it. An earlier, factorial version of the optimiser also searched
along scaled Cauchy and Gauss--Newton directions for any remaining decrease
of the objective (Appendix~\ref{app:nova}); the production solver, used for
the August release and for all current results, does not. The first start
begins at the white-light depth with a flat spectral shape and the theory
limb darkening; the second adds a fixed 250~ppm perturbation to the shape and
small offsets to the limb darkening. For the pilot recoveries, a third start
is placed at the midpoint of the first two, a rule which uses no information
from the injected spectrum. Among the strict starts, the one with the lowest
value of the objective is selected. The policy files still list agreement
tolerances between starts of 0.1 and 0.5~ppm on the two summary grids, but
the current closure does not test them; the agreement is reported instead.

\paragraph{White-light fit.}
Besides the wide geometry priors, $k_o$ is normal with mean 0.1255 and
standard deviation 0.085, $u_1$ and $u_2$ are each uniform on $[0,1]$, the
baseline level and slope are normal about 1 and 0 with standard deviations
0.05 and 0.01, and the noise multiplier is normal with mean 1.4 and standard
deviation 2.0, all copied from the public Ahsoka white-light configuration of
\citet{LouieEtAl2025}. dynesty runs with 2{,}048 live points, multi-ellipsoid
bounds and random-walk proposals, and stops when the estimated remaining
contribution to the evidence, $\Delta\ln Z$, falls below 0.01. The three
failed emcee chains reached only about 20 to 27 autocorrelation lengths.

\paragraph{Grazing-mass rule.}
The central 95\% interval of the grazing mass is estimated from 200 bootstrap
resamplings of the nested-sampling run, made with dynesty's own resampling,
and its half-width must be at most 0.01. If no sample with positive weight is
grazing, the upper end of the interval is set to
$\min(1,\,3/\lfloor N_{\mathrm{ESS}}\rfloor)$, where $N_{\mathrm{ESS}}$ is
the Kish effective sample size; this ``rule of three'' is only an
approximation, since nested samples are not independent. A fit fails if the
upper end is 5\% or more, if the sampler did not reach its stopping
criterion, or if the interval is too wide. In the test on saved posteriors,
the three failed emcee chains had grazing masses of about 12\%. The impact parameter passed to the spectral fit is computed from the weighted
medians of $a/R_\star$ and $i$. Diagnostics of multimodality and of the
posterior tails are reported, but play no part in acceptance.

\paragraph{Uncertainty ensemble.}
Each block of 64 donor integrations starts at a randomly chosen
out-of-transit integration and continues within the same out-of-transit
interval, $[0,230)$ or $[600,720)$, wrapping round to the start of that
interval if necessary; blocks are placed one after another until all 720
integrations are filled. A block of 128 or 192 integrations drawn from the
120-integration interval $[600,720)$ therefore repeats some of its frames.
Each realisation keeps the real ERR and data-quality arrays, and any sample
which becomes invalid in the draw is also flagged. A realisation whose
white-light fit fails or whose starts are not strict is recorded as a
failure; in the two block-length sets it stays in the registered denominator
of 32. The 147 reporting bins are finer than the native basis functions
below 1~$\mu\mathrm{m}$ and above 2~$\mu\mathrm{m}$, which is why the
reporting map has rank 128. A held-out error counts as covered if its
statistic lies below the 95th percentile of the $F$ distribution, and the
validation fails if a one-sided exact binomial test rejects 95\% coverage at
a significance level of 0.025, which happens when 27 or fewer of the 32 are
covered. The Mardia skewness and kurtosis of the 256 estimation errors are
compared with those of 512 simulated Gaussian ensembles of the same size, and
a departure at the same significance level is reported as a warning. The
validation is run once all 352 realisations have finished.

\subsection{The injector}

\paragraph{Baseline.}
The baseline is program 1353, observation 101, visit 1, integrations 0 to 228,
stored in two raw segments of 178 and 51 integrations. Each integration has
eight nondestructive reads of the $256\times2048$ subarray, with one frame per
group, a group time of 5.494~s and no dropped frames, so the accumulated
exposure times run from 5.494 to 43.952~s. The real transit and its buffer integrations occupy integrations 230 to 599.
In the rate-level control, the count change of the final read is divided by
its accumulated exposure time of 43.952~s and added to the calibrated rate
images of the null without the flat-field factor, since those images are
already flat-fielded; the uncertainty and quality arrays of the null are left
untouched.

\paragraph{Source zones.}
The core is a projection of the stellar spectrum through the public ATOCA
spectral-trace model with a wing correction, a per-column flux normalisation
and a per-pixel out-of-transit response calibrated on the same data; this
projection was made earlier in the project and is reused exactly as it was.
Outside the core, the measured light is estimated in bands of 64 columns and
separately for odd and even rows, and the nonnegativity constraint preserves
the signed total of each band over the whole region outside the core. The
zones contain 119{,}277 core pixels and 257{,}554 wing and 121{,}811 far
pixels eligible for the empirical estimate. Far pixels with no measured
light (33{,}758) receive nothing. Every pixel outside the core which is not
eligible for the empirical estimate keeps the modelled target light of the
core projection: these are 10{,}838 pixels flagged in the F277W exposure, a direct image used
to locate contaminating field sources (of
the 12{,}808 flagged pixels, 1{,}820 lie in the core and 150 are reference
pixels) and 4{,}600 pixels without a finite out-of-transit measurement. The
10{,}208 reference pixels receive no light. The core and outer components
are projected separately and summed, and the sum must reproduce the source
map to a relative precision of $5\times10^{-13}$ (an absolute precision of
$10^{-10}\,\mathrm{DN\,s^{-1}}$). In the core, the stellar spectrum is held
at its endpoint value outside its calibration interval, which only affects
the reads of integration~0; the outer component varies linearly between
endpoints defined at the start and the end of the full interval covered by
the reads.

\paragraph{Measured far-light member.}
The far-light dip is measured in the 102{,}308 far pixels which lie inside
the pipeline's background mask, comparing 315 in-transit with 350
out-of-transit integrations and excluding 55 buffer integrations; averaged
over all of these pixels, the signal drops by 0.834\% during the real
transit. The dip of each column is divided by the mean fractional decrement
of the accepted real white-light model over the same 315 integrations,
averaged over the two orders (0.013629), and the same achromatic value is
used for every column; a later refit of the white-light curve changes this
value by only $3\times10^{-7}$, which leaves the member unchanged. The ratio is averaged
over a 128-column window, using only the 1{,}884 columns with a direct
measurement, and bounded to $[0,1]$. The resulting multiplier has a minimum
of 0.103, a median of 0.606 and a maximum of 1, with 84 columns at the upper
bound, and this member injects 61.5\% of the measured far light in total.

\paragraph{Transit tables.}
The injected spectrum is defined on 15{,}326 wavelength nodes. The TRIP
routine uses 20{,}000 annuli; the injector calls the public ExoTETHyS code for
the stellar intensity profile and the total stellar flux, and replaces only
the occulted-flux calculation with a version which evaluates the same
expression one planet position at a time, and which must reproduce the
public routine bit for bit in a regression test before every run. The
MPS-ATLAS intensities are those for $T_{\mathrm{eff}}=6550$~K, $\log g=4.149$
and $[\mathrm{M}/\mathrm{H}]=-0.25$. Grids of 41 and 61 Gauss--Legendre nodes
must agree to better than 1~ppm, and the 61-node result is used; the same
read window is used for every pixel. The per-read tables are checked against
independently computed tables of the transit averaged over whole
integrations: the final read window of each integration must reproduce the
integration average to within $10^{-6}$~ppm, and the integration-averaged
tables must be unocculted in the first and last integrations to within
$5\times10^{-13}$. Integration~0 must carry no transit signal at all, i.e.\ a
remaining fraction of exactly 1 and a time moment of exactly zero. In the
three tables used in the pilot, integration~0 carried a round-off error of
$2\times10^{-16}$ in this fraction, which was verified to be harmless and set
to exactly 1 with explicit approval; in all new tables, any such departure
is recorded as a failure.

\paragraph{Transport.}
The bridge from raw counts to the input of the linearity step passes through
the pipeline's data-quality initialisation, saturation, superbias and
reference-pixel steps; its offsets are computed once from a pair of raw and
pre-linearity frames of the baseline. The bridge must reproduce the input of
the linearity step to within 0.015625~DN, the delivered change must match
the requested change to within the same tolerance, and the linearity
polynomial must be inverted to within $10^{-8}$~DN. The two orders are summed
before transport. Samples flagged as saturated in the group quality array,
or not finite at the linearity input, keep their original values, and the
part of the decrement which could not be delivered is recorded. The
effective flat equals one at pixels flagged in \texttt{jwst\_niriss\_flat\_0275},
whose flags are left in place, and on the original two-order support of
507{,}606 pixels it ranges from 0.2021 to 1.5027. The modified ramps are written as floating-point values. No scene mask is used anywhere in the
injection. In NOVA's own detector processing, the group-level $1/f$
correction is applied between the reference-pixel correction and the
linearity step, so the counts which its linearity step sees differ from those
used in the inversion by this correction.

\paragraph{Detector processing.}
NOVA's chain runs the JWST pipeline \citep{JWSTPipeline2025} up to and
including the reference-pixel correction, then the group-level background and
$1/f$ steps of exoTEDRF 2.4.3 \citep{Radica2024}, then the pipeline's
linearity, jump-detection, ramp-fitting and gain steps, without a dark-current
subtraction, and finally Stage~2 of exoTEDRF, which subtracts the scaled
background model, applies the flat field and corrects bad pixels with
thresholds of 15 in space and time. The group-level background step
subtracts the scaled background only to estimate the $1/f$ noise and adds it
back afterwards. On the injected inputs, the background and $1/f$ steps take
the first and last 50 integrations as their out-of-transit reference. The mask built by the $1/f$ step at run time must be identical to the
reference F277W mask of 12{,}808 pixels. No scene mask of the project's own,
i.e.\ a veto of whole detector columns affected by field stars, is
applied. exoTEDRF 2.4.3 is the version with which
NOVA's input has always been prepared; the comparison pipeline is the later
release 2.5.0.

\paragraph{Delivery checks.}
The written files must have the expected number of reads, the count changes
must be nonpositive, and the protected pixels (the reference pixels and every
pixel outside the support of the injected light) must be bit-for-bit
identical to the input, as must all other extensions of the raw files and
every saturated or inactive sample. The aperture sums of each order and
column are recorded before and after the $1/f$ step, and the frames written
by the step must reproduce the recorded sums exactly. In the grey test, the
gain $g_m$ of each pair of order and column is the regression of the
aperture signal after the $1/f$ step on that before it over integrations 50
to 178, each signal being the difference between the injected run and the
null. The departure from the secant, $16{,}000\,[g_{0.5}-(g_0+g_1)/2]$~ppm, is
evaluated in 2{,}188 pairs of order and column (1{,}742 in order 1 between 1.0
and 2.7~$\mu\mathrm{m}$ and 446 in order 2 between 0.65 and
0.85~$\mu\mathrm{m}$), in which all three gains are finite and the signal
before the correction exceeds 1~$\mathrm{DN\,s^{-1}}$.

\paragraph{Pilot atmospheres and per-visit checks.}
The reference atmosphere assumes rainout condensation, a temperature of
2600~K, a surface gravity of $5\,\mathrm{m\,s^{-2}}$, 200 times solar
metallicity, $\mathrm{C/O}=0.56$, a haze factor of 1100 and a cloud factor of
1. The second atmosphere assumes local condensation, 2400~K,
$10\,\mathrm{m\,s^{-2}}$, 200 times solar metallicity, $\mathrm{C/O}=1.0$, a
haze factor of 10 and a cloud factor of 0.2; it was chosen, before any score
had been seen, as the first distinct case of the existing list of
atmospheres. For each further visit, the core map must not exceed the
observed out-of-transit aperture light in any of five regions (order 1 at
1.0--1.8, 1.8--2.3 and 2.3--2.7~$\mu\mathrm{m}$, and order 2 at 0.65--0.85
and 0.85--1.0~$\mu\mathrm{m}$), and the wing map must differ between two folds
of the data by at most 0.15\% in source fraction, with a held-out wing
residual within 0.25\%.

\paragraph{Comparison pipelines.}
Ahsoka fits each order with the \texttt{emcee} sampler in Eureka!, using 200
walkers and 1{,}000 steps, of which the first 500 are discarded, and fixes
$t_0$, $a/R_\star$ and $i$ in the spectral fit of each order to the posterior
medians of that order's white-light fit. exoTEDRF fits a joint two-order
white-light curve with its own nested sampler.
\texttt{transitspectroscopy} traces the median image with its
double-Gaussian template, extracts a box aperture with a radius of 15 pixels,
fits its joint two-order white-light curve by nested sampling with
\texttt{juliet} and a Mat\'ern-3/2 process for each order, fits the spectral
light curves with ExoTETHyS limb darkening in Kipping coordinates, and uses
ten
spline knots, a 5-sample, $5\sigma$ temporal outlier filter and top-hat
passbands. The published Ahsoka priors put the mid-transit time about 2.5~h
after the end of the injected window, with $a/R_\star\approx7$.

\paragraph{Reporting.}
The reporting grid has 150 bins at $R=100$ between 0.63 and
2.81~$\mu\mathrm{m}$. Bins 147 and 148 lie beyond the validated range of the
public wavelength calibration, and bin 149 is excluded by construction, as in
the earlier release, so the last retained bin ends at $2.740\,\mu\mathrm{m}$.
NOVA's map requires every retained bin to be completely covered; for the
comparison pipelines, the coverage of each bin by their native channels is
reported. The six contiguous reporting ranges have boundaries at 1.0, 1.6,
1.8, 2.3 and 2.7~$\mu\mathrm{m}$, and the two summary ranges are below
2.1~$\mu\mathrm{m}$ (120 bins) and from 2.12 to 2.76~$\mu\mathrm{m}$ inclusive
(26 bins); the one bin whose centre lies between 2.10 and
2.12~$\mu\mathrm{m}$ belongs to neither.
